\centerline{\bf \Huge Тренировочная олимпиада}
\centerline{\bf \Huge Муницип 8 класс}

\begin{flushright}
        \scriptsize{Каждая задача оценивается в 4 балла.}
\end{flushright}

\problem{8.1} Найдите наибольшее натуральное число $n$, при котором (7!)! делится на $n$ !.
Здесь $k$ ! означает произведение $1 \cdot 2 \cdot 3 \cdot \ldots \cdot k$.

\problem{8.2} Саша и Маша пришли в магазин. Саша купил 3 пакета молока, 7 пачек творога и 5 йогуртов. Маша купила 2 пакета молока, 10 пачек творога и 6 йогуртов. Саша потратил на всё 980 рублей, а Маша потратила 1160 рублей. Сколько суммарно стоит один пакет молока, одна пачка творога и один йогурт?

\problem{8.3} Внутри трапеции $A B C D$ с основаниями $A D$ и $B C$ отметили точку $O$. Оказалось, что $A O=B O=C O=B C$ и $D A=D O=D C$. Сколько градусов составляет угол $B A O$ ?

\problem{8.4} В компанию «Рожки и Лапки» устроилось некоторое количество тружеников и 310 лентяев. Каждый день на работу приходило по 50 человек, причём в конце рабочего дня 25 из них говорили: «Сегодня на работу пришло ровно 20 лентяев». Известно, что труженики всегда говорили правду, а лентяи всегда лгали. Через $n$ дней оказалось, что каждый из лентяев сходил на работу ровно 1 раз.

(a) \textit{(2 балла)} Найдите наибольшее возможное значение $n$.

(б) \textit{(2 балла)} Найдите наименьшее возможное значение $n$.

\problem{8.5} У Пети есть три палочки, длины которых равны $a$ см, $b$ см и $c$ см. Известно, что числа $a, b$ и $c$ натуральные, различные, и $a c=b^2$. Петя смог сложить из своих палочек треугольник.

(a) \textit{(2 балла)} Какой наименьший периметр мог быть у этого треугольника?

(б) \textit{(2 балла)} Пусть $b=72$. Какой периметр мог быть у этого треугольника? Укажите все варианты в любом порядке.